\documentclass[10pt]{article}

% ---- NeurIPS-like single-column layout (self-contained; no external conference .sty) ----
\usepackage[letterpaper,left=1.5in,right=1.5in,top=1in,bottom=1in]{geometry}
\usepackage{mathptmx}            % Times text + math (NeurIPS look)
\usepackage[T1]{fontenc}
\usepackage{amsmath,amssymb}
\usepackage{graphicx}
\usepackage{booktabs}
\usepackage{array}
\usepackage{xcolor}
\usepackage[numbers]{natbib}
\usepackage{url}
\usepackage[hidelinks]{hyperref}

\setlength{\parindent}{0pt}
\setlength{\parskip}{0.5em}

% NeurIPS-style abstract block: centered heading between two rules, narrower margins
\renewenvironment{abstract}{%
  \vspace{0.5em}\begin{center}\rule{\textwidth}{0.5pt}\end{center}\vspace{-0.3em}
  {\centering\textbf{Abstract}\par}\vspace{0.3em}
  \begin{quote}\small}{%
  \end{quote}\vspace{-0.3em}\begin{center}\rule{\textwidth}{0.5pt}\end{center}}

\title{\textbf{Over-Squashing {\itshape vs.} Adversarial Robustness in GNNs:\\
An ``Add {\itshape vs.} Remove Edge'' Investigation}}
\author{
  Supanut Kompayak (st126055) \quad Dechathon Niamsa-Ard (st126235)\\
  Deep Learning Course Project, Asian Institute of Technology\\
  \texttt{\{st126055, st126235\}@ait.asia}
}
\date{Midway Report --- \today}

\begin{document}
\maketitle

\begin{abstract}
Graph neural networks (GNNs) suffer from two seemingly unrelated structural problems:
\emph{over-squashing}, where long-range information is compressed through bottlenecks and lost,
and \emph{adversarial fragility}, where a few maliciously edited edges flip predictions. Both
concern how sensitive a node's output is to structural change elsewhere. As a learning project, we
\emph{investigate} whether methods that mitigate over-squashing also help adversarial robustness,
through the lens of \textbf{adding vs.\ removing edges}: over-squashing rewiring methods either
\emph{add} edges (FoSR) or \emph{remove} them (spectral pruning), while adversarial attacks add
edges and adversarial defenses (GCN-Jaccard) remove them. On Cora node classification we compare a
vanilla GCN, an add-based rewiring (FoSR), a remove-based spectral pruning (Jamadandi ProxyDelete),
and a purpose-built defense (GCN-Jaccard) under the Metattack poisoning attack, reporting clean
accuracy, robust accuracy, and the robustness gap, plus a Jacobian-based over-squashing probe. We
report only real, reproducible numbers; negative results are acceptable. Preliminary work establishes
a GCN/Cora baseline of $80.30 \pm 0.47\%$ and validates the full attack--defense pipeline end-to-end.
\end{abstract}

\section{Introduction}
Message-passing GNNs propagate information hop-by-hop along edges. Two known failure modes both stem
from how information flows through graph structure. \textbf{Over-squashing}~\citep{alon2021bottleneck,
topping2022understanding,digiovanni2023oversquashing}: reaching distant nodes requires many layers, but
the exponentially growing long-range signal is squeezed through structural bottlenecks into a
fixed-size vector and lost --- the output is \emph{not sensitive enough} to distant signal.
\textbf{Adversarial fragility}~\citep{zugner2019metattack,zugner2018nettack}: an attacker adds or
removes a few edges and message passing pulls in the bad signal, flipping the prediction --- the output
is \emph{too sensitive} to malicious structural change. Both are about the sensitivity of a node's
output to structural change elsewhere; this conceptual bridge is what we investigate, not a claim we
prove.

The sharper framing we use is \textbf{add vs.\ remove edge}. Over-squashing rewiring mostly
\emph{adds} edges to raise the spectral gap and relieve bottlenecks (FoSR~\citep{karhadkar2023fosr},
SDRF~\citep{topping2022understanding}); adversarial attacks typically \emph{add} edges between
dissimilar nodes; adversarial defenses \emph{remove} suspicious edges (GCN-Jaccard~\citep{wu2019jaccard},
Pro-GNN~\citep{jin2020prognn}). Thus adding edges to fix over-squashing might \emph{hurt} robustness
(more edges to exploit), while a \emph{remove}-based over-squashing method --- spectral
pruning~\citep{jamadandi2024spectral} --- might transfer to robustness because it removes edges like a
defense does. ``Add vs.\ remove: which helps robustness, and why?'' is the spine of this project.

\section{Research Questions}
\begin{itemize}
  \item \textbf{RQ2 (lead).} Do over-squashing mitigations (add-based FoSR vs.\ remove-based spectral
  pruning) improve adversarial robustness relative to (a) an undefended GCN and (b) a purpose-built
  defense (GCN-Jaccard), and at what cost to clean accuracy?
  \item \textbf{RQ1 (secondary).} Does structural perturbation measurably worsen an over-squashing
  measure? Is there a relationship between perturbation strength and long-range information loss?
  \item \textbf{RQ3 (stretch).} Can a single measure (Jacobian sensitivity) capture both
  over-squashing severity and adversarial fragility on the same scale?
\end{itemize}

\section{Data}
Cora is the workhorse; everything else is optional. All statistics are confirmed at load time.

\begin{table}[h]
\centering
\small
\begin{tabular}{@{}llll@{}}
\toprule
\textbf{Dataset} & \textbf{Task} & \textbf{Size} & \textbf{Role} \\
\midrule
Cora (Planetoid)     & node clf.            & 2{,}708 nodes / 5{,}278 undir.\ edges / 7 cls / 1{,}433 feat & \textbf{primary} \\
Citeseer (Planetoid) & node clf.            & 3{,}327 / 4{,}552 / 6 / 3{,}703                              & 2nd setting \\
Peptides-func (LRGB) & graph clf.\ (AP)     & 15{,}535 graphs, $\sim$150 avg nodes                          & stretch \\
ZINC (subset)        & graph reg.\ (MAE)    & 12{,}000 graphs, $\sim$23 avg nodes                          & stretch \\
\bottomrule
\end{tabular}
\caption{Datasets. Cora is the primary node-level and standard adversarial-attack testbed.
Long-range graph-classification sets (LRGB, ZINC) are stretch goals for the over-squashing probe.}
\end{table}

\section{Methods and Approach}
\paragraph{Backbone.} GCN~\citep{kipf2017gcn} is the primary backbone (2 layers, hidden $16$, dropout
$0.5$). GAT and GraphSAGE are stretch.

\paragraph{The method table (the design).} Each method occupies a position on the add/remove axis:

\begin{table}[h]
\centering
\small
\begin{tabular}{@{}llll@{}}
\toprule
\textbf{Method} & \textbf{Targets} & \textbf{Edge op} & \textbf{Role} \\
\midrule
GCN (vanilla)               & ---                 & none   & baseline \\
FoSR                        & over-squashing      & \textbf{add}    & the ``add'' arm \\
Spectral pruning (ProxyDelete) & over-squashing   & \textbf{remove} & the ``remove'' arm \\
GCN-Jaccard                 & adversarial defense & \textbf{remove} & purpose-built defense reference \\
DropEdge                    & regularization      & remove (random) & random-removal sanity baseline \\
\bottomrule
\end{tabular}
\caption{What each comparison teaches: \emph{FoSR vs.\ ProxyDelete} isolates the add/remove
\emph{direction} among over-squashing fixes; \emph{ProxyDelete vs.\ GCN-Jaccard} asks whether
spectral-criterion removal matches similarity-criterion removal; \emph{all vs.\ vanilla} asks whether
any of this helps robustness, and at what clean-accuracy cost.}
\end{table}

\paragraph{Attack and defense.} The attack is Metattack~\citep{zugner2019metattack}, a global
meta-learning poisoning attack and the standard testbed for robust-accuracy benchmarks, at perturbation
rates $\{0.05, 0.10\}$ (with $0.20$ as stretch). The defense reference is
GCN-Jaccard~\citep{wu2019jaccard}, which prunes edges between dissimilar nodes. We use
DeepRobust~\citep{li2021deeprobust} for both. The remove-based over-squashing arm is the ProxyDelete
spectral pruning of \citet{jamadandi2024spectral}, which deletes edges to maximize the spectral gap;
the add-based arm is FoSR~\citep{karhadkar2023fosr}, adapted from graph- to node-classification.

\paragraph{Pipeline order (the operational spine).} Metattack is a \emph{poisoning} attack, so the
order of operations is part of the design, not an implementation detail:
\[
\text{clean graph} \xrightarrow{\text{Metattack}} \text{poisoned } G'
\xrightarrow{\;M\;} G'' \xrightarrow{\text{train GCN}} \text{eval},
\quad M \in \{\text{none, FoSR, ProxyDelete, GCN-Jaccard}\}.
\]
The attack precedes training by definition of poisoning; the method $M$ is applied to the received
(possibly poisoned) graph before training, mirroring deployment and matching how DeepRobust evaluates
GCN-Jaccard. Because a method changes the graph even with no attack, its clean-accuracy cost must be
measured separately, giving four graph conditions per method: clean; clean${+}M$; poisoned; and
poisoned${+}M$. The headline table reports, per method, clean accuracy, robust accuracy at ptb
$\{0.05, 0.10\}$, and the robustness gap (clean $-$ robust).

\paragraph{Metrics.} (i) \emph{Clean accuracy}; (ii) \emph{robust accuracy} under Metattack;
(iii) \emph{robustness gap} $=$ clean $-$ robust (headline); and (iv) an \emph{over-squashing measure}
based on Jacobian sensitivity $\partial \mathbf{h}_v / \partial \mathbf{x}_u$ for distant node pairs
$(u,v)$~\citep{digiovanni2023oversquashing}, computed via autograd (scoped to a few far pairs first,
then expanded). The focus of the investigation is RQ2; RQ1 follows almost for free once the probe
exists.

\section{Discussion}
The grid directly answers RQ2. The attacker mostly \emph{adds} dissimilar edges; GCN-Jaccard
\emph{removes} exactly those (which is why removal defends); FoSR \emph{adds} more edges (does adding
hurt robustness?); ProxyDelete \emph{removes} edges by a spectral criterion (does over-squashing-motivated
removal happen to remove attacker edges too?). \emph{FoSR vs.\ ProxyDelete} isolates the effect of edge
direction among methods with the same over-squashing goal, while \emph{ProxyDelete vs.\ GCN-Jaccard}
isolates whether the removal \emph{criterion} matters --- i.e.\ whether generic spectral removal suffices
or removal must target the attack. Our working hypothesis is that add-based rewiring is neutral-to-harmful
for robustness while remove-based pruning partially transfers, but a negative or weak result is an
acceptable and informative outcome. The Jacobian probe links the two phenomena: if perturbation increases
long-range sensitivity (RQ1) and the same measure tracks fragility (RQ3), the conceptual bridge is
empirically supported on one scale.

\paragraph{Caveat.} We use a non-adaptive attacker: Metattack attacks a vanilla GCN surrogate and is
unaware of the defense applied afterward (the DeepRobust standard). This flatters defenses relative to an
adaptive attacker; we note it as a limitation and treat adaptive attacks as a stretch goal.

\section{Summary of Progress}
\paragraph{Done.} (1) A reproducible GCN/Cora baseline of $\mathbf{80.30 \pm 0.47\%}$ test accuracy
(5 seeds, CPU), confirming Cora statistics at load. (2) The remove-arm code (Jamadandi ProxyDelete) was
verified to exist and be usable (\texttt{RelationalML/SpectralPruningBraess}; pure NetworkX/SciPy,
CPU-friendly, runs on Cora node classification), upgrading the remove-arm from a random fallback to a
principled spectral method. (3) The DeepRobust attack--defense pipeline was installed (working around a
broken \texttt{gensim} pin via \texttt{--no-deps} plus prebuilt PyG extension wheels) and validated
end-to-end on CPU; Table~\ref{tab:spike} reports the smoke-test numbers.

\begin{table}[h]
\centering
\small
\begin{tabular}{@{}lccc@{}}
\toprule
\textbf{Setting (Cora, Metattack ptb $=0.05$)} & \textbf{Clean} & \textbf{Robust (poisoned)} & \textbf{Gap} \\
\midrule
GCN (Planetoid public split, 5 seeds)  & $80.30 \pm 0.47$ & --- & --- \\
GCN (DeepRobust attack split)          & $84.05$ & $77.72$ & $6.34$ \\
GCN-Jaccard (DeepRobust attack split)  & ---     & $79.38$ & --- \\
\bottomrule
\end{tabular}
\caption{Preliminary numbers (\%). Top row: the clean baseline on the standard Planetoid split. Bottom
rows: the Task~1 end-to-end pipeline smoke test on the attack split used by robustness papers (single
seed, untuned) --- enough to confirm the path works; the full tuned grid is Task~2.}
\label{tab:spike}
\end{table}

\paragraph{To do.} Build the core pipeline and headline table over $\{$clean, ptb $0.05$, ptb $0.10\}$
(Task~2); adapt FoSR to node tasks (Task~3); implement the scoped Jacobian probe (Task~4); add Citeseer
and integrate ProxyDelete (Task~5); prepare slides and the progress summary (Task~6).

\paragraph{Obstacles and concerns.} (i) DeepRobust ($0.2.9$, 2023) is brittle against current
PyTorch/NumPy --- resolved by a clean dependency install and by running on CPU for these small graphs,
sidestepping Blackwell-GPU (sm\_120) issues. (ii) FoSR's reference code targets graph classification and
must be adapted to node tasks. (iii) The Jacobian probe is the hardest component and is scoped small
first. (iv) The timeline is tight, so Cora is prioritized and all other datasets and the adaptive-attack
setting are explicitly stretch goals.

\paragraph{Reproducibility.} All numbers in this report are produced by fixed-seed scripts (the
GCN/Cora baseline and the DeepRobust attack--defense pipeline spike) and are reproducible end-to-end;
we report only real, measured results, and the codebase is organized so that each table can be
regenerated from a single command.

\bibliographystyle{plainnat}
\bibliography{references}

\end{document}
